\documentclass[11pt,a4paper]{article}
\usepackage[utf8]{inputenc}
\usepackage[T1]{fontenc}
\usepackage{amsfonts}
\usepackage[french]{babel}
\frenchbsetup{StandardLists=true} % à inclure si on utilise \usepackage[francais]{babel}
\usepackage{enumitem}
\usepackage{amssymb}
\usepackage{amsmath}
\usepackage[left=2cm,right=2cm,top=2cm,bottom=2cm]{geometry}
\usepackage{multirow}
\usepackage[dvipsnames]{xcolor}

\usepackage[T1]{fontenc}
\usepackage{fouriernc}
\usepackage{graphicx}

\usepackage{setspace}
\singlespacing

\usepackage{pstricks}
\usepackage{xkeyval}
\usepackage{pst-labo}


\usepackage{multicol}
\usepackage{caption}
\usepackage{fancyhdr}
\pagestyle{fancy}
\renewcommand\headrulewidth{1pt}
\fancyhead[L]{Lycée Denis Diderot }
\fancyhead[R]{Classe de 1 Bac Pro }
\setcounter{page}{1}\fancyfoot[C]{page \thepage}
\fancyfoot[R]{Année 2024-2025}
\fancyfoot[L]{Alexis RUIZ}
\usepackage{multirow,tabularx,booktabs}
\usepackage[tikz]{bclogo} 

\newcommand{\cadreblanc}[1]{%
\medskip\framebox[\linewidth][c]{%
\rule{0mm}{#1mm}}\par
\medskip}

\usepackage{awesomebox}
\setlength{\aweboxvskip}{0mm}
\usepackage{pifont}
\usepackage{chemmacros}
\chemsetup{modules=all}
\setlength{\columnseprule}{.25mm}

\renewcommand{\arraystretch}{2.5}
\newcommand*\acocher{\quitvmode{\fboxsep0pt \fboxrule0.8pt \fbox{\vrule height1.8ex depth.1ex width0pt \vrule height0pt depth0pt width1.9ex }}}

\newcolumntype{C}[1]{>{\centering\arraybackslash }b{#1}}

\newcommand{\code}{\awesomebox[black]{2pt}{

\includegraphics[scale=0.15]{Chapitre 1 - Energie et puissance électrique/images/frame.png}
}{black}}

\newcommand{\moteur}{\awesomebox[black]{2pt}{

\includegraphics[scale=1]{images/moteur.jpg} 
}{black}}

\newcommand{\defconclusion}{\awesomebox[black]{2pt}{\rotatebox{90}{\Large{\textbf{Conclusion}}}}{black}}

\newcommand{\defbox}{\awesomebox[black]{2pt}{
\includegraphics[scale=0.5]{Chapitre 1 - Energie et puissance électrique/images/appel exam.png}}{black}}

  \newenvironment{itemize*}%  % listes compactes
    {\begin{itemize}%
      \setlength{\itemsep}{0pt}%
      \setlength{\parskip}{0pt}}%
    {\end{itemize}}

\frenchbsetup{StandardLists=true}



\begin{document}

\flushleft \textbf{NOM : \hfill Classe : \hspace{3cm} ~~\\[2mm]
Prénom:\\[0,4cm]
}
\hrule\
\\[0,2cm]

\begin{center}
 \LARGE{TP1  - Comment mesurer la puissance moyenne consommée par un appareil électrique en régime alternatif ?   }
\end{center}

\vspace{3 mm}

\normalsize

\hrule\
\\[0,2cm]
   \begin{center}
   \textbf{Le soin et la rédaction interviendront dans la notation.}
   \end{center} 

\noindent%
\begin{minipage}{.48\textwidth}%

    
 \begin{bclogo}[couleur=blue!30 ,  arrondi =0.1 , logo = \bcinfo, epBarre =2]{ Matériel }

\begin{itemize*}
    \item 1 résistor de 33 $\mathrm{\Omega}$
    \item 1 bobine avec L = 70 mH et r = 26,5 $\mathrm{\Omega}$
    \item 1 wattmètre
    \item 1 générateur 6V - 50 Hz
    \item 1 plaque de test 
    \item 5 fils de connexion
\end{itemize*}
\end{bclogo}

\end{minipage}%
\hfill
\begin{minipage}{.45\textwidth}%
\includegraphics[width=\textwidth]{Chapitre 3 - Puissance consommée par un appareil éléctrique/images/image.png}
\end{minipage}%



\fcolorbox{black}{blue!30}{\textbf{Protocole}}
\begin{itemize*}
    \item Réaliser le montage schématisé ci-dessus.
    \defbox{\fbox{
\begin{minipage}{8cm}{

Appeler M Ruiz pour qu'il vérifie le montage.}

\end{minipage}
}}

    \item En utilisant les fonctionnalités du wattmètre, mesurez la puissance électrique consommée $P$, la tension $U$ et l'intensité du courant $I$. Écrire les valeurs dans le tableau ci-dessous. 

    \item Dans le montage, ajoutez la bobine en série avec le résistor. 

    \defbox{\fbox{
\begin{minipage}{8cm}{

Appeler M Ruiz pour qu'il vérifie le montage.}

\end{minipage}
}}

    \item Mesurer à nouveau $P$ , $U$ et $I$.


\fcolorbox{black}{blue!30}{\textbf{Mesures}}\\
\ding{202} Complétez le tableau suivant : 

\begin{center}
    

    \begin{tabular}{|C{2cm}|C{2cm}||C{2cm}|C{2cm}||C{2cm}|C{2cm}|}
         \hline
         Dipole & P (en W) &  U (en V)  & I (en A) & U $\times I$ & $\dfrac{P}{U \times I}$\\
         \hline
         Résistor &  &   &  &  &  \\
         \hline
         Bobine  &  &  &  &  & \\
         \hline
 
    \end{tabular}
\end{center}




\fcolorbox{black}{blue!30}{\textbf{Interprétation }} \\[5mm]
\ding{203} Proposer pour le résistor une relation entre $P$ , $U$ et $I$. \\

\cadreblanc{15}
\vspace{3mm}

\ding{204} Cette relation est-elle vérifiée pour une bobine ? 
\vspace{3mm}
\dotfill 





\fcolorbox{black}{blue!30}{\textbf{Validation}}\\[5mm]
Avec un compteur d'énergie portatif, on a mesuré la puissance électrique consommée $P$, dite "active", la tension efficace $U$ et $I$ l'intensité efficace pour un four composé de résistor (image 1) et pour un moteur contenant des bobine (image 2) \\[5mm]

\ding{205} Observez les valeurs mesurées et choisissez l'écran qui correspond à chaque appareil. \\[3mm]

\begin{center}
    Justifier votre choix. 
\end{center}

\noindent%
\begin{minipage}{.45\textwidth}%
\begin{center}
\includegraphics[scale=0.4]{Chapitre 3 - Puissance consommée par un appareil éléctrique/images/ecran 1.png}

\cadreblanc{20}
\end{center}
\end{minipage}%
\hfill
\begin{minipage}{.45\textwidth}%
\begin{center}
\includegraphics[scale=0.4]{Chapitre 3 - Puissance consommée par un appareil éléctrique/images/ecran 2.png}

\cadreblanc{20}
\end{center}
\end{minipage}%









\end{itemize*}


\end{document}